% Arquivo: latex/secao4_discussao_resultados_DRAFT.tex
% Capítulo 4 - Discussão dos Resultados
% DRAFT para revisão - NÃO incluir diretamente em main.tex sem aprovação

\section{DISCUSSÃO DOS RESULTADOS}
\label{sec:results_syst_review}

A literatura empírica sobre os instrumentos da Política Nacional de Desenvolvimento Regional apresenta expressiva assimetria em termos de volume de evidências: os Fundos Constitucionais concentram a quase totalidade dos estudos, enquanto os Fundos de Desenvolvimento e os Incentivos Fiscais constituem vertentes consideravelmente menos desenvolvidas. Os resultados obtidos variam substancialmente conforme a especificação econométrica adotada, o período amostral, a escala geográfica de análise e o controle de características locais não observáveis. Esta seção apresenta os principais achados da literatura, organizados por instrumento e abordagem metodológica.

\subsection{Avaliações de Impacto dos Fundos Constitucionais sobre o PIB}
\label{subsec:fc-pib}

Os Fundos Constitucionais são os instrumentos da PNDR mais amplamente avaliados pela literatura empírica: dentre os 35 estudos aprovados nesta revisão, 29 investigam os efeitos desses fundos sobre variáveis de resultado relacionadas ao produto interno bruto ou ao mercado de trabalho. O FNE é o fundo com maior número de avaliações (24 estudos), seguido pelo FCO (11 estudos) e pelo FNO (8 estudos). A predominância dos Fundos Constitucionais na literatura decorre de sua maior longevidade em relação aos demais instrumentos da PNDR -- com operação contínua desde 1989 --, da disponibilidade de dados municipais anuais de PIB a partir de 2002, do elevado volume de recursos envolvidos e da capilaridade dos programas em todos os municípios das macrorregiões-alvo. A quase totalidade dos estudos adota o PIB \textit{per capita} municipal como variável de resultado, tratando-o como aproximação da renda média local, dada a escassez de alternativas em frequência anual e cobertura municipal. Os resultados obtidos pela literatura variam consideravelmente conforme a especificação econométrica adotada, o período amostral, a escala geográfica de análise e o controle de características locais não observáveis. As subseções seguintes organizam os achados por abordagem metodológica, permitindo identificar padrões gerais e divergências nas estimativas de impacto.

\subsubsection{Modelos de Efeitos Fixos}

Os primeiros estudos a avaliar os efeitos dos Fundos Constitucionais sobre o PIB municipal utilizaram modelos de painel com efeitos fixos de município e tempo. \citeonline{Resende2014a} analisa os efeitos do FNE sobre o crescimento do PIB \textit{per capita} dos municípios nordestinos no período 2004-2010, utilizando médias em três subperíodos (2004-2006, 2006-2008 e 2008-2010) para minimizar viés de ciclos de negócios. Os resultados apontam efeito positivo do FNE-total em nível municipal (0,021) e microrregional (0,032), mas não significativo em nível mesorregional, o que os autores atribuem à maior heterogeneidade em regiões geográficas maiores. As estimativas sugerem que o efeito positivo estaria concentrado nos empréstimos ao setor agropecuário, porém os resultados não se mostraram robustos à inclusão de efeitos fixos temporais. Para o FNO, aplicando a mesma especificação econométrica à região Norte, o autor obtém relação inversa entre o fundo e o crescimento municipal, com ausência de robustez nas análises setoriais quando se controla por efeitos fixos temporais. \citeonline{OliveiraResendeGoncalves2017} aplicam especificações de efeitos fixos juntamente com outros métodos (PSM e GPS) para avaliar o FCO no estado de Goiás, obtendo evidências de efeito positivo sobre o crescimento municipal, com magnitude variando conforme o método empregado. Mais recentemente, \citeonline{SilvaFilhoetal2024} adotam painel de efeitos fixos para avaliar conjuntamente os três fundos constitucionais, não obtendo evidências consistentes de associação positiva com o crescimento do PIB \textit{per capita} para o período 2016-2019, resultado que contrasta com os estudos anteriores e pode estar associado ao curto intervalo temporal analisado ou a mudanças no contexto macroeconômico.

Os autores desses estudos ressaltam limitações metodológicas importantes. A principal dificuldade consiste na identificação causal dos efeitos, dado que a alocação dos fundos é endógena à demanda por crédito subsidiado, o que tende a privilegiar municípios com maior capacidade de absorção e infraestrutura socioeconômica mais desenvolvida. Adicionalmente, apontam-se dificuldades em mensurar os custos fiscais efetivos dos programas, a possibilidade de existência de peso morto (investimentos que seriam realizados mesmo sem subsídio) e potenciais efeitos de deslocamento de fatores produtivos entre municípios vizinhos. A ausência de robustez dos resultados à inclusão de controles temporais sugere que parte do efeito estimado pode capturar tendências macroeconômicas comuns aos municípios tratados e não tratados, em vez do efeito causal dos fundos. Em síntese, os modelos de efeitos fixos indicam correlação positiva dos fundos constitucionais com o crescimento econômico local, com efeitos mais pronunciados em nível municipal e maior magnitude para o FCO, mas a fragilidade da identificação causal e a ausência de robustez em diversas especificações limitam a interpretação causal dos resultados.

\subsubsection{Modelos de Painel Espacial}

Uma vertente da literatura admite a possibilidade de interdependência econômica entre unidades espaciais geograficamente próximas, o que pode ocasionar efeitos indiretos de transbordamento da política para outras regiões, expandindo a área potencialmente beneficiada. Por outro lado, o efeito poderia consistir em vazamentos de recursos das regiões-alvo da política em direção a regiões de maior dinamismo econômico. Portanto, entender a natureza da dinâmica espacial é importante para um planejamento mais eficiente da política. \citeonline{CravoResende2015} utilizam modelo de painel SDM (\textit{Spatial Durbin Model}) e estimam efeito significativo do FCO sobre o PIB \textit{per capita}, com evidência de que uma elevação de 1 ponto percentual (p.p.) na relação FCO/PIB está associada a um incremento de cerca de 0,09 p.p. na taxa média anual de crescimento do PIB \textit{per capita} municipal nos dois anos seguintes. Na análise setorial, o resultado seria atribuído à modalidade de crédito FCO-Empresarial, em que o efeito seria na ordem de 0,13 p.p. de incremento. Porém, o resultado é não significativo para o FNE e negativo para o FNO. Com respeito aos efeitos indiretos (transbordamento espacial), obteve-se resultado relevante, mas não positivo, para o FCO, o que pode estar ligado à migração de fatores de produção.

\citeonline{OliveiraLimaArriel2016} usam modelo espacial para avaliar efeitos do FCO nos municípios de Goiás e concluem que efeitos diretos são menores quando se admite a ocorrência de efeitos espaciais, sugerindo que modelos não espaciais tendem a superestimar o impacto real dos fundos. Entretanto, as evidências obtidas por \citeonline{ResendeCravoPires2014} e \citeonline{CravoResende2015} apontam em direção contrária, com os efeitos diretos do FCO maiores em modelos espaciais do que em modelos não espaciais. Essa divergência pode estar associada a diferenças na cobertura geográfica (Goiás versus toda a região Centro-Oeste), no tipo de matriz de pesos espaciais empregada (\textit{queen} de contiguidade versus distância inversa) e no período de análise considerado. Os resultados de \citeonline{OliveiraLimaArriel2016} indicam que cada 1\% de aumento no valor dos financiamentos está associado a aproximadamente 0,1\% de crescimento no PIB \textit{per capita} do município, usando especificação SAR com matriz espacial \textit{queen} do tipo contiguidade e dados de 2004 a 2011.

\citeonline{ResendeSilvaFilho2017} distinguem-se dos trabalhos anteriores por incorporar na análise a heterogeneidade definida pela própria PNDR por meio das tipologias regionais. Em modelo não espacial de efeitos fixos, o efeito estimado do FNE é significativo para municípios pertencentes às tipologias dinâmica e baixa renda. Quando se admite dependência espacial, há evidências de transbordamento espacial dos investimentos do FNE quando o município diretamente beneficiado tem tipologia dinâmica. As estimativas apontam que para cada 1 p.p. de aumento na relação FNE/PIB em municípios dinâmicos há um incremento de 0,07 p.p. na taxa média anual de crescimento do PIB \textit{per capita} do próprio município nos três anos seguintes e de 0,33 p.p. na taxa média dos municípios vizinhos. Para os autores, esses resultados ressaltam a relevância de se considerar a dimensão espacial nas análises e também o caráter heterogêneo dos efeitos do FNE nos territórios. \citeonline{ResendeSilvaFilho2018} adotam a mesma abordagem econométrica, mas incluindo também as regiões Norte e Centro-Oeste, a fim de avaliar conjuntamente os fundos FNE, FNO e FCO. Os resultados são positivos para efeitos diretos e indiretos em municípios de tipologia dinâmica quando se usam valores acumulados na variável dos fundos (dois primeiros anos de cada subperíodo). Os resultados apontam que um aumento de 1 p.p. na relação FC/PIB em municípios dinâmicos está associado a 0,03 p.p. de incremento no crescimento médio do PIB \textit{per capita} do próprio município pelos três anos seguintes e 0,12 p.p. de incremento nos municípios vizinhos. \citeonline{SilvaFilhoetal2023} não obtêm evidências de associação espacial positiva entre os financiamentos do FNE, FNO e FCO nos anos de 2016 a 2018 e o crescimento médio do PIB \textit{per capita} dos municípios de 2016 a 2019, porém, o resultado é positivo sobre o PIB \textit{per capita} de 2019. Os autores concluem que os efeitos dos fundos são em sua maior parte esporádicos e pontuais.

Em resumo, os estudos convergem para a relevância de se incorporar a dimensão espacial nas análises de efeito dos fundos constitucionais, justificada pela ocorrência de correlação espacial significativa nas variáveis envolvidas. Porém, há menor consenso quando se trata da magnitude dos efeitos. Em geral, o resultado estaria condicionado a outras características econômicas locais. O efeito direto mostrou-se significativo em municípios de tipologia dinâmica e microrregiões de baixa renda, com incrementos de 0,03 a 0,14 p.p. no crescimento anual do PIB \textit{per capita} municipal por até dois anos, para cada 1 p.p. de aumento na relação Fundo/PIB. Para o efeito indireto, as estimativas variam entre 0,1 e 0,4 p.p. de incremento no crescimento para cada 1 p.p. de aumento em Fundo/PIB, também concentrado em municípios de tipologia dinâmica e microrregiões de baixa renda. Em análise mais desagregada, evidências apontam para associação entre o FCO e taxa de crescimento, atribuída à linha de crédito empresarial. No FNE, os resultados estariam ligados ao setor industrial e de serviços. Já para o FNO não há nenhuma evidência positiva e significativa. No aspecto metodológico, nota-se a necessidade de estudos mais robustos, que combinem intervalo de tempo mais longo, especificações espaciais alternativas, correção de endogeneidade e maior número de controles (inclusive variáveis de confusão da própria PNDR).

\subsubsection{Modelos Não Lineares e de Eficiência}

\citeonline{Linharesetal2014} usam modelo \textit{threshold} e obtêm evidências de que o efeito do FNE é pouco significativo para municípios com PIB \textit{per capita} inicial abaixo de R\$ 2.143 ou acima de R\$ 7.406. Entre esses limites, os autores identificam dois grupos distintos. O primeiro é formado por municípios de PIB \textit{per capita} inicial entre R\$ 2.143 e R\$ 3.866. Nesse grupo, estima-se que 1 p.p. de aumento médio em FNE/PIB no período 2002-2006 esteja associado a um incremento de 0,08 p.p. no crescimento médio do PIB \textit{per capita} entre 2002 e 2008. No segundo grupo, de municípios com PIB \textit{per capita} inicial entre R\$ 3.866 e R\$ 7.406, o efeito estimado foi maior, na ordem de 0,11 p.p. Os resultados sugerem que o FNE tem maior efetividade em municípios com níveis intermediários de renda, reafirmando os achados de \citeonline{ResendeSilvaFilho2017} e \citeonline{ResendeSilvaFilho2018}. Em municípios de baixa renda, porém, os resultados não foram significativos, o que indica a necessidade de aprimoramento dos mecanismos da política para melhor adequação às características e necessidades dessas regiões.

\citeonline{IrffiAraujoBastos2016} usam modelagem quantílica para estimar a significância e heterogeneidade de efeitos do FNE ao longo da distribuição de crescimento dos municípios, empregando estimação em dois estágios e a localização de agências bancárias do BNB como instrumento para a alocação de recursos do fundo. Os efeitos foram positivos e significativos para o FNE total e para as linhas de crédito associadas à agricultura e pecuária. No entanto, os largos intervalos de confiança não permitiram constatar presença de heterogeneidade ao longo da curva de crescimento. \citeonline{Carneiro2018} analisa a relação entre a participação do FNE e a eficiência média dos municípios nordestinos. Usando modelo de análise envoltória de dados e de fronteira estocástica, o autor identifica elevada eficiência nos municípios pertencentes à região do MATOPIBA, associada à agricultura extensiva voltada à exportação. Adicionalmente, municípios do Semiárido indicaram menor eficiência no crescimento do produto médio, enquanto municípios menos densos e mais distantes das respectivas capitais e com empresas menores tendem a fazer melhor uso dos financiamentos.

\subsubsection{Modelos Dinâmicos}

\citeonline{Cambota2019} adotam modelo de painel dinâmico e instrumentos sintéticos para corrigir o viés de endogeneidade na variável PIB defasada e na variável de interesse FNE/PIB, usando dados de 2003 a 2014. Os autores confirmam a importância desse tipo de controle para a obtenção de estimativas mais realistas do impacto da política. As estimativas apontam associação positiva entre o FNE e o crescimento econômico dos municípios, na ordem de 2,96 p.p. para cada 1 p.p. de aumento em FNE/PIB no ano anterior. Cabe notar que essa magnitude é consideravelmente superior à dos demais estudos revisados, cujos coeficientes situam-se tipicamente entre 0,03 e 0,13 p.p. A discrepância pode ser atribuída à diferente especificação do modelo (painel dinâmico com instrumentos sintéticos), ao período mais longo de análise (2003-2014) ou a possível sobreestimação decorrente da proliferação de instrumentos no procedimento GMM, aspecto reconhecido na literatura como potencial fonte de viés.

\subsubsection{Modelos de Equilíbrio Geral}

Em \citeonline{NascimentoHaddad2017}, um cenário hipotético de retirada do FNE e realocação dos recursos em gastos correntes do governo resultaria a longo prazo em uma redução líquida de 0,16\% no PIB nacional e em maior concentração regional da produção. O estudo também sugere que os fundos têm papel mais relevante no crescimento de longo prazo, via elevação do estoque de capital e da capacidade produtiva local, frente à alternativa apontada. No curto prazo, ambos os cenários seriam equivalentes em termos de ativação da demanda. \citeonline{Ribeiroetal2020} indicam que os investimentos do FNE realizados em 2014 e 2015 teriam elevado o PIB da região Nordeste em 3,51\% até o ano 2025. \citeonline{GoncalvesBragaGurgel2022} constroem dois cenários alternativos envolvendo o término da política de subsídios, e ambos indicam como consequência a redução na atividade econômica na região Nordeste, puxada principalmente pela redução do produto da atividade agropecuária.

\subsubsection{Modelos Quase Experimentais}

\citeonline{Resende2014c} usa o método de primeira diferença para avaliar o efeito de empréstimos industriais do FNE sobre o crescimento do PIB \textit{per capita} municipal. Usando taxas de crescimento do PIB \textit{per capita} nos períodos 2000-2003 e 2000-2006, o autor não encontra efeito significativo das operações envolvendo o FNE-Indústria durante o ano 2000, seja em nível municipal ou microrregional. Argumenta-se que o resultado pode estar associado à elevada participação do setor agropecuário nas operações do FNE, setor em que há predomínio da informalidade, o que dificulta sua captura pelas medições do PIB.

\citeonline{OliveiraSilveiraNeto2020} adotam desenho quase-experimental que combina regressão com descontinuidade geográfica (RDD) e modelo de diferenças em diferenças. Os autores constroem um painel de dados decenal (1980, 1990, 2000 e 2010) composto por municípios dos estados de Minas Gerais e Espírito Santo que ingressaram nos critérios de elegibilidade ao FNE em períodos recentes (dados pré e pós-ingresso) e por municípios vizinhos imediatos que nunca tenham sido tratados pelo FNE. Os resultados obtidos levantam dúvidas sobre a efetividade do FNE, sugerindo que o Fundo não teve influência sobre as variáveis de produto e renda nos municípios avaliados no ano 2010. O estudo ressalta a importância de melhor vinculação do crédito ao aumento da produtividade, mediante condições diferenciadas e combinação com outros instrumentos.

Por outro lado, \citeonline{Carneiroetal2024a} obtêm evidências de efeito positivo e significativo do FNE sobre o PIB \textit{per capita} municipal. Os autores usam dados de 2002 a 2019 junto a estimadores DID dois estágios e DID escalonado, combinando três instrumentos de política regional: FNE, FDNE e incentivo fiscal da SUDENE. O estudo aponta efeito médio de tratamento de 8,8\% a 11\% no PIB \textit{per capita} do município beneficiado em relação ao grupo de controle, sendo atribuído principalmente ao setor serviços e administração pública. O estudo de eventos aponta que o efeito de tratamento é crescente e inicia no quinto ano após o primeiro tratamento. Quanto à sinergia com os demais instrumentos, o resultado é não significativo. \citeonline{MonteIrffiBastosCarneiro2025} exploram heterogeneidades de gênero, tipo de tomador e magnitude dos empréstimos (efeito dosagem) do FNE, aplicando o método \textit{Generalized Propensity Score} (GPS). O estudo encontra correlação positiva entre a participação do FNE-Mulheres no valor total do programa e a variação do PIB \textit{per capita} municipal, o que sinaliza uma eficácia maior do programa nesse grupo populacional. Efeito semelhante é obtido em relação à variável participação de tomadores pessoa jurídica no total de financiamentos. Com respeito ao efeito dosagem, os resultados apontam efeitos positivos no PIB \textit{per capita} somente na última e mais elevada dose de tratamento.

Em conjunto, as evidências sobre o efeito dos Fundos Constitucionais no PIB \textit{per capita} revelam um quadro heterogêneo e condicionado à especificação metodológica adotada. Os modelos de efeitos fixos e espaciais indicam associação positiva para o FCO e resultados ambíguos para FNE e FNO, enquanto os modelos quase experimentais apontam tanto para ausência de efeito quanto para impactos positivos significativos. Os modelos não lineares e de eficiência sugerem que a eficácia dos fundos é condicionada ao nível inicial de renda do município, com efeitos mais pronunciados em faixas intermediárias de desenvolvimento. As simulações de equilíbrio geral reafirmam a importância dos fundos para a atividade econômica regional em cenários contrafactuais de retirada da política. A constatação de que os efeitos são mais pronunciados em municípios com níveis intermediários de renda, e nulos ou inconclusivos em municípios de baixa renda, pode ser interpretada à luz da literatura sobre capacidade de absorção e armadilhas de pobreza: municípios de baixa renda podem carecer do capital humano e da infraestrutura básica necessários para converter o crédito subsidiado em crescimento econômico sustentado; municípios de renda intermediária possuiriam o mínimo de capacidade necessária para ativar o efeito multiplicador dos investimentos; e municípios de alta renda estariam próximos do efeito marginal decrescente do capital.

\subsection{Avaliações de Impacto dos Fundos Constitucionais sobre o Mercado de Trabalho}

Dezesseis estudos aprovados nesta revisão avaliam o impacto dos Fundos Constitucionais sobre variáveis do mercado de trabalho, com foco predominante no emprego formal, massa salarial e salário médio. \citeonline{SilvaResendeSilveiraNeto2009} foi o primeiro estudo a investigar o impacto dos empréstimos dos FCs sobre a contratação de trabalhadores nas empresas beneficiadas. O estudo faz a interseção dos registros de empresas beneficiadas com os dados da RAIS (Relação Anual de Informações Sociais), obtendo a quantidade de vínculos de empregos formais de 449 empresas (NE=211, N=174, CO=74) nos anos de 2000 e 2003. Aplicando o método \textit{Propensity Score Matching (PSM)}, os autores concluem que nas empresas beneficiadas pelo FNE a quantidade de empregos cresceu em média 61,2 pontos percentuais a mais do que nas empresas não beneficiadas. No entanto, os valores não foram significativos para o salário médio. E ainda, para o FNO e o FCO, não foi obtido efeito positivo em nenhum dos indicadores avaliados.

\citeonline{SoaresSousaNeto2009} também aplicam o método PSM, mas fazem uso de um conjunto mais rico de informações sobre o FNE, o que permitiu ampliar consideravelmente o número de empresas e o horizonte de tempo investigado. Além disso, os autores fazem uma análise ano a ano dos efeitos do fundo, o que permite estimativas mais eficientes, com quantidades maiores de empresas nos primeiros anos pós-tratamento. Com dados do período de 2000 a 2006, no horizonte de 3 anos pós-tratamento há 1.104 empresas e em 5 anos há 314. Os resultados obtidos reforçam as evidências obtidas por \citeonline{SilvaResendeSilveiraNeto2009}. No primeiro ano pós-tratamento o crescimento do estoque de empregos é em média 7,96 pontos percentuais maior nas empresas beneficiadas, relativamente às demais. Além disso, a trajetória é crescente nos anos subsequentes, alcançando 33,59 pontos percentuais no terceiro ano e 132,23 no quinto ano. Se confirmados, esses resultados atribuem grande relevância para a política regional, indicando que o FNE tem papel importante na geração de empregos formais nas empresas beneficiadas. Os autores também destacam que o efeito de maior monta estaria associado às empresas financiadas entre 1999 e 2001, as únicas no horizonte de 5 anos pós-tratamento, estando mais concentradas no setor industrial. O mesmo estudo ainda analisa possíveis efeitos do FNE sobre a massa salarial, obtendo resultados significativos e com dinâmica semelhante ao do efeito no emprego. Como consequência, a variável salário médio não apresentou variação significativa como decorrência do FNE, indicando que as novas contratações teriam ocorrido ao mesmo nível dos salários vigentes.

\citeonline{OliveiraMenezesResende2018a} apresentam como inovação principal o uso de variável de tratamento do tipo contínua, permitindo analisar possíveis heterogeneidades no efeito do fundo decorrentes da intensidade do tratamento (valor do empréstimo), em contraste à abordagem dicotômica de tratamento presente nos estudos anteriores. Os autores usam o método \textit{Generalized Propensity Score Matching (GPS)} e dados de financiamentos a empresas do estado de Goiás realizados pelo FCO por meio da linha Empresarial (associada à indústria, comércio e serviços) no período de 2004 a 2011. Os resultados indicam relação positiva e significativa com o crescimento do emprego apenas em um intervalo específico da distribuição de tratamento e limitada ao período 2004-2008 (pré-crise), cenário semelhante ao obtido na abordagem PSM com variável de tratamento dicotômica. Com respeito ao salário médio, não se obteve nenhum resultado significativo. \citeonline{OliveiraMenezesResende2018b} expandem a abordagem de \citeonline{OliveiraMenezesResende2018a} para os demais fundos constitucionais. O dado mais relevante obtido indica efeito decrescente dos fundos sobre as variáveis de interesse, o que significa que quanto maior o valor do empréstimo, menor o efeito sobre o crescimento do emprego e do salário. Para os autores, isso se justifica do ponto de vista teórico pela propriedade de retornos marginais decrescentes da função de produção neoclássica.

\citeonline{CunhaJuniorSoares2024} contribuem com a literatura ao usar o método de função dose-resposta para captar efeitos heterogêneos e correção de viés de seleção por meio de variável instrumental. Para garantir comparabilidade, os autores optam por usar a mesma base de dados de \citeonline{SoaresSousaNeto2009}. Na variável crescimento do emprego, o modelo IV aponta efeito médio de 357,54\%, acima de 133,94\% do modelo OLS. Para o crescimento da massa salarial, o valor médio é 220,96\% no modelo IV contra 188,55\% em OLS. Em ambas as variáveis a restrição de exogeneidade leva a subestimação do efeito médio e superestimação do ponto ótimo de tratamento, sinalizando para efeitos mais significativos em pontos mais baixos da distribuição de tratamento, quando a endogeneidade é corrigida.

\citeonline{Resende2014c} contribui para uma limitação importante dos métodos PSM e GSM: a incapacidade de controlar por efeitos não observáveis que possam influenciar tanto a probabilidade de tratamento quanto seu resultado (apenas características observáveis). Usando abordagem de primeira diferença e dados de 1999 a 2006, o autor identifica efeito positivo do FNE-Indústria sobre o crescimento do número de empregos formais nas empresas beneficiadas, com magnitude entre 10,83 e 16,11 pontos percentuais ao ano. No horizonte de três anos pós-tratamento, o valor é de 30 a 56 pontos percentuais a mais, de certa forma consistente com os resultados de \citeonline{SilvaResendeSilveiraNeto2009} e \citeonline{SoaresSousaNeto2009}.

\citeonline{OliveiraLimaArriel2016} é o primeiro estudo a avaliar o efeito dos fundos constitucionais sobre variáveis agregadas do mercado de trabalho. O estudo adota os municípios de Goiás como unidade de análise e aplica o método de painel espacial autorregressivo (SAR), a fim de controlar também por características da vizinhança do município. Os resultados indicam efeitos positivos sobre emprego e salário, porém, sem diferenças relevantes em relação ao modelo de painel clássico com efeitos fixos, o que pode ser decorrência de baixa robustez dos resultados. \citeonline{RiegerLimaRodrigues2020} trazem duas contribuições principais: análise dos efeitos agregados do FNE por setor de atividade econômica e o uso de modelo de painel dinâmico para amenizar possível viés de endogeneidade. Os resultados apontam relação positiva entre os investimentos do FNE e o nível de emprego no município, com maior efeito sobre o setor primário (1,1\% de aumento no emprego para 1\% de aumento no investimento do FNE), seguido por terciário (0,53\%) e secundário (0,13\%). Esses resultados são observados para os investimentos correntes do FNE. Para o FNE defasado em um ano, o valor é significativo somente para o setor secundário e em menor magnitude (0,06\%). Sugere-se assim que o FNE tem efeito positivo e imediato sobre o emprego, especialmente no setor primário, setor mais intensivo em mão de obra.

Em maior nível de agregação, \citeonline{Ribeiroetal2020} aplicam um modelo dinâmico interregional de equilíbrio geral computável (CGE) para estimar os efeitos dos investimentos FNE sobre o emprego total da região Nordeste, dos estados e dos setores de atividade. As evidências apontam que os investimentos do FNE nos anos 2014 e 2015 estão associados a uma variação acumulada de 2,75\% na quantidade de empregos na região Nordeste até o ano de 2025, valor pouco menor que o estimado para a variação do PIB. Na análise por estados, Piauí, Ceará e Rio Grande do Norte também apresentam a maior variação acumulada do emprego, entre 3,19\% e 4,31\%, e Maranhão e Bahia têm as menores variações, de 0,45\% a 2,16\%. No aspecto setorial, também não há perdedores, com agropecuária (8\%) e construção (5,3\%) os setores com maior variação do emprego.

\citeonline{DanielBraga2020} empregam o método diferenças-em-diferenças, que controla pela heterogeneidade não observável das empresas e possibilita uma estimativa do efeito causal do fundo. Além disso, têm como diferencial a análise dos efeitos microeconômicos do FNO por porte da empresa beneficiada, tipo de crédito concedido e horizonte de tempo. No geral, empresas beneficiadas apresentaram desempenho superior às não beneficiadas, em curto e médio prazo, em termos de geração de emprego e renda. No entanto, médias e grandes empresas beneficiadas por financiamento a capital de giro e custeio apresentam desempenho similar ao grupo de controle, no que diz respeito à geração de emprego e renda. Além disso, os autores apontam que o impacto positivo do crédito com finalidade capital de giro e custeio é em média maior no curto prazo do que no longo prazo. Por outro lado, em empresas beneficiadas pela finalidade investimento, os impactos são mais observáveis a longo prazo. Os efeitos assimétricos indicam que a finalidade investimento promove a acumulação de capital, o que aumenta a produtividade do trabalho a longo prazo, pois o impacto é positivo para o trabalho e a renda (salário médio).

\citeonline{OliveiraSilveiraNeto2021} utilizam modelo de descontinuidade geográfica (RDD) para estimar o impacto causal do FNE sobre o estoque de emprego, distinguindo-se dos estudos anteriores pela robustez metodológica e por analisar impacto sobre o emprego formal e informal, decorrência do uso de dados censitários. Os resultados indicam que o acesso ao FNE gerou impacto positivo e significativo sobre o estoque de empregos formais e informais no ano 2010, com magnitude 5,4\% (agropecuária), 2,8\% (indústria) e 3,5\% (serviços). Porém, impacto sobre a taxa de informalidade foi nulo.

Em síntese, as avaliações sobre o mercado de trabalho apresentam maior convergência do que as relativas ao PIB: há evidências consistentes de efeito positivo do FNE sobre a geração de emprego formal nas empresas beneficiadas, com magnitude crescente ao longo dos primeiros anos pós-tratamento. Os estudos também convergem em apontar efeitos mais pronunciados em micro e pequenas empresas e na modalidade de crédito para investimento. Por outro lado, a variável salário médio não apresenta variação significativa na maioria das especificações, sugerindo que os novos postos de trabalho são criados ao nível salarial vigente, sem ganhos adicionais de produtividade. Os resultados para FNO e FCO são mais limitados e inconclusivos, refletindo menor cobertura na literatura.

\subsection{Avaliações de Impacto dos Fundos de Desenvolvimento}

Os Fundos de Desenvolvimento (FDNE, FDA e FDCO) constituem instrumentos mais recentes da PNDR, com operação sistemática a partir de 2010-2013, o que explica a menor quantidade de estudos empíricos disponíveis. Dentre os 35 estudos aprovados nesta revisão, 8 avaliam algum aspecto dos Fundos de Desenvolvimento. Os estudos recentes de \citeonline{Braz2024}, \citeonline{CarneiroCostaIrffi2024b} e \citeonline{Irffietal2025} são as primeiras evidências empíricas disponíveis de efeito de Fundos de Desenvolvimento no PIB local. \citeonline{Braz2024} usam modelo DID escalonado e estimam que municípios que receberam empreendimentos apoiados pelo FDNE experimentam, até 2021, uma elevação média de 24\% no PIB \textit{per capita} em relação ao grupo de controle. \citeonline{CarneiroCostaIrffi2024b}, ao avaliarem simultaneamente os efeitos do FNE e dos incentivos fiscais da SUDENE, não encontraram significância estatística para os incentivos fiscais, mas identificaram impacto positivo para o fundo FDNE, com elevação de cerca de 23\% no PIB \textit{per capita} municipal para uma elevação de 10\% nos recursos do FDNE. Complementando, \citeonline{Irffietal2025} estimam que a instalação de parques de geração de energia eólica gera aumento médio de 21\% no PIB \textit{per capita} municipal. O estudo avalia todos os parques eólicos da região Nordeste. Embora o tamanho amostral não permita avaliar especificamente os parques eólicos financiados pelo FDNE, a proximidade com os valores obtidos nos estudos anteriores reforça a consistência das conclusões. Estas primeiras evidências sugerem que os Fundos de Desenvolvimento podem ter impacto positivo expressivo sobre o PIB local, com magnitudes superiores às estimadas para os Fundos Constitucionais. Todavia, a base empírica ainda é restrita, concentrada em três estudos recentes que empregam metodologias semelhantes (DiD escalonado) e analisam períodos parcialmente sobrepostos, o que limita a generalização dos resultados. Avaliações futuras sobre outros resultados -- como emprego, renda domiciliar, desigualdade e pobreza -- constituem prioridade para a agenda de pesquisa sobre os Fundos de Desenvolvimento.

\subsection{Avaliações de Impacto dos Incentivos Fiscais}

A avaliação empírica dos incentivos fiscais da SUDENE e SUDAM constitui a vertente menos desenvolvida da literatura sobre a PNDR. Dentre os 35 estudos selecionados, 9 avaliam explicitamente estes instrumentos, refletindo a operação mais recente dos programas (concessões sistemáticas a partir de 2010) e as dificuldades de acesso a dados detalhados de renúncia fiscal em nível de empresa. Essa escassez de evidências representa uma das lacunas mais relevantes identificadas nesta revisão. No âmbito do mercado de trabalho, \citeonline{Garsousetal2017} estimam aumento de 30\% na quantidade de empregos no setor turístico como consequência dos incentivos fiscais da SUDENE, usando modelo de diferenças-em-diferenças. \citeonline{Braz2023} avaliam o efeito sobre os vínculos totais da economia e obtêm evidências de elevação em 3,2\%. \citeonline{CostaCarneiroIrffi2024} analisam os efeitos dos incentivos fiscais da SUDENE sobre a criação de empregos em empresas da região Nordeste, empregando o método \textit{Generalized Synthetic Control}. As evidências indicam efeitos benéficos sobre o aumento do emprego nas empresas beneficiadas, com impactos positivos concentrados no setor manufatureiro intensivo em mão de obra. \citeonline{CarneiroCostaIrffi2024b}, ao incluírem incentivos fiscais conjuntamente com FNE e FDNE em estimador DID escalonado, não encontram significância estatística para os incentivos fiscais sobre o PIB \textit{per capita}. Na literatura de efeitos espaciais, \citeonline{FerreiraIrffiCarneiro2024} apontam a ocorrência de autocorrelação espacial positiva na distribuição de incentivos fiscais da SUDENE. Usando um modelo SAC estático, obtêm efeito indireto significativo do incentivo fiscal da SUDENE sobre o Índice de Desenvolvimento Municipal (IDM), as componentes IDM-saúde e IDM-renda, mas não no IDM-geral, e sem evidências do FDNE. Quanto aos efeitos diretos, o impacto direto do incentivo fiscal sobre o IDM-geral é significativo.

Em síntese, as evidências disponíveis sobre incentivos fiscais são esparsas e inconclusivas, insuficientes para caracterizar de forma robusta seus efeitos sobre o desenvolvimento regional. Nota-se ausência completa de avaliações dos incentivos fiscais da SUDAM, bem como de estudos que empreguem métodos quase experimentais com maior poder de identificação causal para o PIB \textit{per capita}. A ampliação desta literatura constitui prioridade para a agenda de pesquisa sobre a PNDR.

\subsection{Síntese Comparativa dos Resultados}

O panorama de evidências mapeado nesta revisão revela expressiva assimetria entre os instrumentos da PNDR, tanto em volume de estudos quanto em robustez das conclusões. Os Fundos Constitucionais concentram 29 dos 35 estudos aprovados, com evidências de efeito positivo sobre emprego (maior convergência entre métodos) e resultados heterogêneos sobre o PIB \textit{per capita} (fortemente condicionados ao nível de renda do município, à especificação metodológica adotada e ao período amostral). Os Fundos de Desenvolvimento, avaliados apenas recentemente por 8 estudos, apresentam indícios de impacto positivo expressivo sobre o PIB local (21-24\%), porém com base empírica ainda restrita a três estudos com metodologia semelhante. Os Incentivos Fiscais permanecem como a vertente menos investigada, com 9 estudos e evidências inconclusivas.

A convergência de resultados varia substancialmente conforme a variável de resultado analisada. Para o mercado de trabalho, há consenso de que o FNE gera impacto positivo sobre o emprego formal nas empresas beneficiadas, com efeitos crescentes ao longo dos primeiros anos pós-tratamento e maior magnitude em micro e pequenas empresas. A variável salário médio, por sua vez, não apresenta variação significativa na maioria das especificações, sugerindo que os novos postos de trabalho são criados sem ganhos adicionais de produtividade. Para o PIB \textit{per capita}, os resultados são mais divergentes: modelos de efeitos fixos e espaciais indicam correlação positiva mas fragilidades de identificação causal; modelos não lineares e de eficiência demonstram que a eficácia dos fundos é condicionada ao nível inicial de renda do município, com efeitos mais pronunciados em faixas intermediárias de desenvolvimento; e modelos quase experimentais apresentam tanto ausência de efeito quanto impactos positivos significativos, a depender da especificação e do período analisado.

Esta heterogeneidade de resultados pode ser interpretada à luz da literatura sobre capacidade de absorção e armadilhas de pobreza: municípios de baixa renda podem carecer do capital humano e da infraestrutura básica necessários para converter crédito subsidiado em crescimento econômico sustentado; municípios de renda intermediária possuiriam o mínimo de capacidade necessária para ativar o efeito multiplicador dos investimentos; e municípios de alta renda estariam próximos do efeito marginal decrescente do capital. Essa configuração indica que a literatura avançou consideravelmente nos últimos anos, mas ainda carece de avaliações que integrem simultaneamente os três instrumentos e empreguem estratégias de identificação mais robustas para o controle de endogeneidade e dependência espacial. A ampliação da literatura sobre Fundos de Desenvolvimento e Incentivos Fiscais, bem como o uso de métodos quase experimentais com maior poder de identificação causal, constituem prioridades para a agenda de pesquisa sobre a PNDR.
